\documentclass[10pt,a4paper]{article}
\usepackage[T1]{fontenc}
\usepackage[utf8]{inputenc}
\usepackage[margin=1.8cm]{geometry}
\usepackage{lmodern,booktabs,array,tabularx,enumitem,xcolor,hyperref,microtype}
\hypersetup{hidelinks}
\setlength{\parindent}{0pt}
\setlength{\parskip}{4pt}
\setlist{itemsep=1pt,topsep=2pt}
\newcolumntype{L}{>{\raggedright\arraybackslash}X}
\newcommand{\keybox}[1]{\fcolorbox{gray}{gray!8}{\parbox{\dimexpr\linewidth-2\fboxsep-2\fboxrule}{#1}}}
\begin{document}
{\Large\bfseries Claim-level coding protocol}\\[2pt]
Full-text coding of engineering-agent studies and title/abstract screening of the recall-audit records. Version 1, 4 November 2026.

\keybox{\textbf{General rules.}
(1) \emph{Independence}: individual studies are not discussed between coders before both codings are complete.
(2) \emph{Blinding}: coders do not consult earlier codings of the same records, and the coding forms contain no codes from the automated comparison coding.
(3) \emph{Source-based decisions}: every decision cites a page of the publication. Evidence that is not reported is treated as absent.}

\section{Object of the assessment}
The coding assesses whether the computational results that a study reports for its \emph{headline claim} would support the release of that claim as an engineering result. Only what the publication and the artifacts it refers to report is considered, and the five obligations are assessed individually. The assessment concerns the reported evidence, not the overall quality of the study or the correctness of its engineering. Evidence that is not reported leads to \texttt{UNR} (unresolved), never to \texttt{REF}.

\section{Procedure for one study}
\begin{enumerate}[leftmargin=*]
\item Read the abstract and conclusions of the full text, then the evaluation and results sections.
\item Decide \textbf{eligibility} (Section 3). For an excluded study, record the exclusion code and proceed to the next study.
\item Record the \textbf{headline claim}: the single main computational claim stated in the abstract or conclusions, quoted with its page.
\item Decide whether \textbf{physical validation is required} (Section 4).
\item Assign the dispositions of \textbf{S, A, N, V, B}, each with a page and a short quotation (Section 5). V is coded \texttt{NR} when it is not required.
\item Record the primary and, where applicable, secondary \textbf{ETF} type (Section 6), the \textbf{evidence tier} (Section 7), and the 16 \textbf{ARCS} items (Section 8).
\item Record the time spent on the study.
\end{enumerate}

\section{Eligibility at full text (stream S1)}
\textbf{Include} if the study describes or evaluates an LLM or agentic system that can formulate, alter, execute, inspect, verify, postprocess, or release artifacts or claims in a computational-engineering or numerical-simulation workflow (rule I-S1), \emph{and} reports some evaluation. Publication period: 1 January 2017 to 15 September 2026. Preprints are eligible.\\
\textbf{Exclude} with the first applicable code: E01 static ML or surrogate model without agent control of simulation artifacts; E05 no material tool use or workflow action (conversational use only); E06 the agent does not act on a numerical engineering simulation artifact or claim; E09 opinion, position, or non-scholarly source without evaluation; E10/E15 the available text is insufficient for a decision (an incorrect or missing full text is noted); E11 duplicate; E16 no substantive publication. Reinforcement-learning agents without a language model do not satisfy I-S1 (code E05). \texttt{UNCERTAIN} is used only when no decision is possible, with an explanatory note.

\section{Requirement of physical validation (V)}
\texttt{V\_required = Y} when the headline claim asserts \emph{any} of the following: the physical performance of a design or product; agreement with physical observations; or a predictive or design decision about a physical system. Otherwise \texttt{V\_required = N}, for example for claims about numerical accuracy with respect to an analytical solution, about software or code generation, or about benchmark pass rates.

\section{Dispositions}
\begin{tabularx}{\textwidth}{>{\bfseries}p{0.04\textwidth} L L L}
\toprule
 & \texttt{SUP} supported & \texttt{REF} refuted & \texttt{UNR} unresolved (default) \\
\midrule
S & The executed model, input, or code is compared with the request (the specification) and matches it for the runs behind the claim. Agreement of the result with a reference solution of the \emph{intended} problem also counts (rule R1). & A reported mismatch between what was run and what was requested, in the evidence behind the claim. & A specification is only declared, or a checking procedure is named without its outcome. \\
\addlinespace
A & The authority of the agent to change the computation was bounded \emph{before} it acted, by an approval gate, a sandbox, or confinement to deterministic typed tools, and the executed effects are stated (R4). & A reported effect outside the permitted scope. & Otherwise. Acceptance by the user \emph{after} the run, or a control that is asserted but not evidenced, is not sufficient. \\
\addlinespace
N & A numerical result for the quantity of interest: the error with respect to an analytical, manufactured, or reference solution; a mesh or time-step refinement study; bounded iterative and discretization errors. For a surrogate, agreement with reference FE or CFD solutions of the same problem (R1). & Reported numerical evidence contradicts the asserted accuracy. & Only execution, compilation, residual reduction, a single mesh, or a statement of convergence. \\
\addlinespace
V & A matched physical comparison or test in the regime of the claim. Possible only if N is \texttt{SUP} (R2). & A reported physical comparison contradicts the claim. & Otherwise, or N not supported. \texttt{NR} if V is not required. \\
\addlinespace
B & Code, data, logs, or models are released (with a link), or retained and explicitly linked to the reported runs (R5). & Results are reported for a configuration that was not run. & Otherwise. Availability on request is not sufficient. \\
\bottomrule
\end{tabularx}

\section{Evidence-transfer failure (ETF)}
The ETF type identifies the substitution on which the headline claim relies. One primary type and, optionally, one secondary type are recorded; \texttt{none} is recorded when the claim remains within its evidence.

\noindent\begin{tabularx}{\textwidth}{>{\bfseries}p{0.05\textwidth} L}
E1 & Execution, compilation, or task completion used as evidence that the intended problem was solved (S). \\
E2 & Execution, convergence, residual, or mesh-quality proxy used as evidence of numerical accuracy (N). \\
E3 & Numerical or benchmark agreement used as evidence of physical validity (V). \\
E4 & Provenance or traceability used as evidence that results are correct or reproducible. \\
E5 & Success on a canonical or benchmark case generalized to a different or general application. \\
E6 & An authorization, guard, or constraint check used as evidence of engineering correctness. \\
\end{tabularx}

\section{Evidence tier}
\textbf{T1}: native numerical or physical evidence for the quantity of interest (error with respect to reference or manufactured solutions, refinement study or GCI, matched experiment). \textbf{T2}: agreement with a benchmark, a reference code, or values from the literature, without an error bound. \textbf{T3}: execution, task-success, or code metrics only. The strongest evidence reported for the headline claim is coded.

\section{ARCS checklist (Y / N / NA)}
\begin{tabularx}{\textwidth}{p{0.03\textwidth} L p{0.03\textwidth} L}
1 & Headline claim states a quantity of interest and an intended use & 9 & Discretization-uncertainty estimate (GCI or equivalent) \\
2 & Approved specification given (geometry, boundary conditions, materials, operating point) & 10 & $y^+$ (CFD) or element quality (FE) reported; NA if neither applies \\
3 & Executed model or input compared with the specification & 11 & Validation data with an uncertainty; NA if V is not required \\
4 & Agent effects (files or settings changed) logged or reported & 12 & Validation conditions match the regime of the claim; NA if V is not required \\
5 & Human approval or authorization gate described & 13 & Failed or discarded agent attempts reported \\
6 & Solver, version, and numerical schemes reported & 14 & Run or artifact identities link results to runs \\
7 & Iterative convergence shown for the quantity of interest, not only for residuals & 15 & Code, prompts, or case files released \\
8 & Mesh or time-step refinement study reported & 16 & Limitations of the claim stated \\
\end{tabularx}

\section{Title/abstract screening}
Decisions are based on the title and abstract; where no abstract is available, the publisher record is consulted through the DOI. \textbf{INCLUDE} if at least one stream criterion holds, with the stream(s) recorded:
\begin{itemize}[leftmargin=*]
\item \textbf{S1} (2017 to 15 September 2026): rule I-S1 above.
\item \textbf{S2} (2017 to 15 September 2026): describes, evaluates, surveys, or formalizes an attack, control, trust mechanism, authorization model, provenance or trace mechanism, or assurance method that is materially applicable to tool-using LLM agents.
\item \textbf{S3} (any date up to 15 September 2026): defines, evaluates, applies, standardizes, or critically reviews a method needed to judge whether a computational model, numerical result, workflow artifact, or simulation-supported claim is credible for an intended use (verification, validation, uncertainty, credibility, provenance, assurance cases).
\end{itemize}
\textbf{EXCLUDE} only when the record clearly fails every stream, with the first applicable code E01--E16 (listed in \texttt{review/search/screening\_reason\_codes\_v4\_1\_1.csv}). \textbf{UNCERTAIN} is used whenever no decision is possible; the record then proceeds to full-text assessment. Exclusion on the title alone is permitted only for an unambiguous scope failure.

\section{Worked examples (hypothetical)}
\textbf{(a)} ``The agent sets up OpenFOAM cases from text and reaches 90\% task success on 20 benchmark cases.'' No comparison with reference solutions; code released. $\Rightarrow$ V\_required N; S \texttt{UNR}, A \texttt{UNR} unless a gate or sandbox is described, N \texttt{UNR}, V \texttt{NR}, B \texttt{SUP}; ETF E1 (secondary E2); T3.\\
\textbf{(b)} ``The agent designs a heat sink that lowers the peak temperature by 12\%.'' Mesh study reported; no experiment. $\Rightarrow$ V\_required Y; N \texttt{SUP}; V \texttt{UNR}; ETF E3; T1.\\
\textbf{(c)} ``Agent-generated FEniCS solvers reach second-order convergence against manufactured solutions; the LLM can only fill a validated JSON schema that is executed by fixed templates.'' $\Rightarrow$ S \texttt{SUP} (R1), A \texttt{SUP} (confinement to typed tools), N \texttt{SUP}; V \texttt{NR}; ETF none; T1.

\section{Calibration and adjudication}
\begin{enumerate}[leftmargin=*]
\item \textbf{Calibration.} Five records are coded jointly before the main coding. Clarifications of the protocol agreed during calibration are recorded, apply to all subsequent coding, and are reported. Calibration records are excluded from agreement statistics.
\item \textbf{Coding.} The remaining records are coded according to this protocol.
\item \textbf{Adjudication.} Disagreements are resolved from the full text, and each decision is recorded with its reason and page reference. The automated comparison coding is consulted only after adjudication.
\end{enumerate}
The procedure applied in the study, in which one author coded all records and a second author verified every record, is reported in Section~2 of the article and in Supplementary Section~S16.2.
\end{document}
